\section{A finite acquired certificate and its intervention meaning}
\label{sec:acquired}

\subsection{A common branch, with no hidden rare-event threshold}
Fix $0<\varepsilon\le1/4$. Conservatively replace a normalized gap
certificate by $\min(g,1)$, and define
\begin{align}
C&=\left\lceil\log_2\frac{128}{\varepsilon\zeta}\right\rceil,
&b&=\lceil C/g\rceil,\\
L&=\left\lceil\frac{1024(H+b+2)}{g\varepsilon}\right\rceil,
&T_{\max}&=\lceil16LC/g\rceil+1.
\label{eq:commonSchedule}
\end{align}
The two models must use these same buffers, block lengths, clock units
and event predicate. Count the learned stationary event probability to a
positive interval $[\ell,u]$, and widen it to
$[p_-,p_+]=[e^{-2E}\ell,e^{2E}u]$ with outward rounding. Refine Count
and require $p_+/p_-\le2$. Under the energy precision below this is a
finite relative-counting task. Both true and learned probabilities then
belong to this interval. All counting failures have their own confidence
budget; this branch rule is used only on the simultaneous good event.

If $p_+\le1/(8L)$, use Theorem~\ref{thm:stability} with $q_0=g/8$.
The accepted structural estimates give $S_Q,S_{\widehat Q}<\varepsilon/32$.
If $p_+>1/(8L)$, then $p_->1/(16L)$. Direct simulation to
$T=\lceil C/(gp_-)\rceil\le T_{\max}$ suffices. For either model,
the killed-semigroup bound
$\Pr_\nu(\tau>T)\le\zeta^{-1/2}e^{-gp_-T}$ is at most
$\varepsilon/128$. Coupling the two finite uniformized trajectories gives
\begin{equation}
\TV(\Law_{\nu,Q}(\mathcal R_{H,s}),
        \Law_{\nu,\widehat Q}(\mathcal R_{H,s}))
\le\sqrt{T\kappa_+/2}+\varepsilon/64.
\label{eq:directStability}
\end{equation}
Thus the direct branch does not reintroduce a numerical $1/p$ cost:
its certified horizon is bounded by $T_{\max}$ before execution.

For example, the following sufficient, deliberately loose precisions
imply a model-law discrepancy below $\varepsilon/4$ in either branch:
\begin{equation}
E\le\frac{\varepsilon\sqrt g}{256},\qquad
a\le\min\left\{1,
\frac{\varepsilon}{256}\sqrt{\frac g{L+1}},
\frac{\varepsilon}{64\sqrt{T_{\max}}}\right\}.
\label{eq:sufficientprecision}
\end{equation}
Indeed $e^a-1\le2a$, and consequently
$D_L\le3\varepsilon^2g/32768$,
$z\le\sqrt3\,\varepsilon/64<\varepsilon/32$,
$d_p\le\varepsilon/128$ and the observed-prefix term is below
$\varepsilon/32$. The geometric term is below $\varepsilon/8$.
The direct bound is smaller. Finite upward rounding can use strict
slack in these inequalities. These are an existence schedule, not an
optimized experimental design and not lower bounds.

\subsection{A certificate that can use an actually measured acceptance}
The uniform schedule is not compulsory. For a chosen $b,L$, suppose
finite calibration provides an outward interval
$\widehat\mu\in[\widehat m_-,\widehat m_+]$ including all numerical and
sampling biases. Equation~\eqref{eq:submass} yields
\begin{equation}
\mu\ge m_-:=
\bigl(\sqrt{\widehat m_-}-\sqrt{D_L}\bigr)_+^2.
\label{eq:observedAcceptance}
\end{equation}
If $m_->0$, use $q_0=\min(m_-,\widehat m_-)$ in the comparison.
For either model its structural bound can be evaluated without knowing
$p$ exactly:
\begin{equation}
S\le\frac{\eta}{1-\eta}
+\frac{1+\eta}{1-\eta}
\frac{H+b+2}{(1+L)m_{\rm low}}
+(1+\eta)p_+(b+1),\qquad \eta=e^{-gb}/\zeta,
\label{eq:instanceStructural}
\end{equation}
where $m_{\rm low}$ is the corresponding acceptance lower bound.
This uses $q(t)\le p(1+t)$ and $q_G=p(1+L)\mu$.
The calibrated proposal must use the intended $G$ and common window;
one cannot insert an unrelated observed acceptance rate. If the
square-root gap closes, this certificate reports ``unresolved'' rather
than a positive floor. Its finite arithmetic and statistical errors are
explicit inputs, not silently erased.

\subsection{Local evidence supplies the uncertainty objects}
Use the accepted reversible, pairwise binary Gibbs acquisition model.
There are $N$ variables, $M$ declared interaction edges, maximum degree
$d$, a finite menu of $K$ constant calibrated interventions $u$, and
full observed model closure. In intervention $u$,
\begin{align}
W^u(x)&=\exp\!\left\{\sum_i(h_i+u_i)x_i+
                         \sum_{ij}J_{ij}x_ix_j\right\},\\
Q^u(x,x^i)&=r_i^u(x_{\partial i})
\sigma\bigl((1-2x_i)(h_i+u_i+\sum_jJ_{ij}x_j)\bigr).
\end{align}
The unknown context clocks have supplied or acquired valid numerical
bounds $0<r_*\le r_i^u(s)\le r^*$. The additive-field model,
local detailed balance, faithful isolated copies, the absence of hidden
memory and the declared graph are premises. Finite observations do not
certify model completeness.

The accepted balanced-field learner, with parameter tolerance $\eta_0$
and kinetic log-rate tolerance $\rho_0$, gives simultaneous envelopes
\begin{equation}
E_u=(N+2M)\eta_0+N(\xi_u+\tau_u),\qquad
a_u=\rho_0+(1+2d)\eta_0+\xi_u+\tau_u.
\label{eq:acquiredEnvelopes}
\end{equation}
Here $\xi_u$ is deployment-field calibration error and $\tau_u$ is
the rational factor/control rounding allowance. Errors in actual timing
and detector outputs must be included in that learner's confidence
certificate; they are not covered by the displayed field allowance.
The original complete acquisition proof and experiment record are
preserved unchanged. In particular, a biased calibration clock is not
assumed to equal an operating clock.

Interval Dobrushin rows, the pinned conditional-probability contrast
condition of the accepted correlated static algorithm, rate bounds and
a mass floor construct common $g,\zeta,R$ for every model in the
retained confidence box. A coarse box may first establish strict
margins; subsequent acquisition refines inside that box. A procedure
with unresolved margins returns that status. It cannot silently turn a
mixing assumption into a learned fact. The static solver is used at
its original graph, contrast and pinned-mixing scope, including its
numerical inverse margins and degree dependence.

Equations~\eqref{eq:acquiredEnvelopes} and
\eqref{eq:sufficientprecision}, or the tighter explicit certificate
\eqref{eq:master}, now specify an acquisition precision sufficient for
the joint event record. This is the new interface. The previous CDF
stability theorem alone supplied none of these marked-path terms.

\subsection{The complete resource bill for the fitted-model route}
Let $D=N+2M$ and $C_{\rm kin}=K\sum_i2^{\deg(i)}$. Omitting only
universal numerical constants, the inherited finite schedules use
\begin{align}
\text{energy reads}\quad&
O\!\left(D\left[\log(2+B_0)+\eta_0^{-2}\right]
                 \log\frac{D(1+\log(2+B_0))}{\delta_E}\right),\\
\text{energy dwell}\quad&
O\!\left(\frac D{r_{\rm cal}}
\left[\log(2+B_0)+\eta_0^{-2}\log(2/\eta_0)\right]
\log\frac{D(1+\log(2+B_0))}{\delta_E}\right),\\
\text{operating releases}\quad&
O\!\left(C_{\rm kin}\rho_0^{-2}
                      \log(C_{\rm kin}/\delta_Q)\right),\\
\text{capped kinetic dwell}\quad&
O\!\left(\frac{C_{\rm kin}}{r_*\rho_0^2}
\log(C_{\rm kin}/\delta_Q)
\log\frac{C_{\rm kin}(1+\rho_0^{-2})}{\delta_Q}\right).
\label{eq:fitCosts}
\end{align}
The original constants, finite censoring and timing construction govern
these summaries. Priced pins and neighbor clamps, context setup, bit
reads, event detector service, field calibration and finite-slew travel
are added at their actual costs. In particular, a pin is not a free
state reset and a detector's clock digits are not a physical response
time guarantee. The degree term and numerical $r_*^{-1}$ can be large.

Once this acquisition is complete, the learned-model sampler pays the
accepted finite static, reverse/forward and clock compilation costs.
Its conservative rare-branch calibration uses
$O(g^{-1}\varepsilon_{\rm sim}^{-2}\log(1/\delta_{\rm sim}))$
length-$L$ Palm proposals; a new draw uses at most
$O(g^{-1}\log(1/\varepsilon_{\rm sim}))$ proposals and a length-$b$
initial segment. Each proposal uses one approximate rare stationary
state and $O(L+\log(1/\theta))$ finite updates, at the prescribed
per-call tolerance $\theta$. Direct-branch updates are bounded by
$T_{\max}$ and the same finite truncation rules. The waiting index is
generated in binary by the preserved finite geometric-clock compiler;
its digits, not its value, are paid. A physical window of duration $h$
requires $H=\lceil Rh\rceil$. Window output and memory scale accordingly.

All formulas depend on local experimental scales, degree, inverse
mixing/static margins, desired accuracy and $\log(1/\zeta)$. There is
no $1/\pi(B)$ or mean-first-event simulation factor. There may still be
exponential costs in degree, physical contrast, mixing margin or
apparatus precision. The earlier thirteen-year serial example is not
made economical by this theorem; its coarse CDF calibration is not a
joint-record certificate.

\subsection{Preparation and decisions that the access really permits}
The abstract comparison is uniform over a common supplied initial law.
For an actual equilibrium-safe mission, use the accepted finite physical
preparation only for menu interventions with the certified bound
$\pi^u(A)\ge1/2$, verified by $p_+^u\le1/2$ from the same widened
count interval. This safe-mass premise is part of the complete operation,
not a consequence of a mixing gap. If this premise is to be certified
only from a two-sided statistical interval, a strict safe-mass slack
(and its required counting precision) is paid; an unresolved boundary
returns that status. In the rare branch the displayed event-mass bound
already supplies ample safe slack. Under the chosen constant intervention,
let the repairable
chain evolve for its certified mixing interval, inspect the state at a
known completion time, retain a safe result, and repeat up to a finite
cap. This requires the physical process to continue through unsafe
states during preparation. It is not an experiment on an irreversibly
destroyed system. Its conditional successful law $\nu_u$ satisfies
\begin{equation}
e^{-\beta_u}\le
\frac{d\nu_u}{d\pi^u(\cdot\mid A)}\le e^{\beta_u},
\qquad A=B^c,
\end{equation}
and its abort probability is at most $\delta_{\rm prep}$. The inherited
inspection, field ramp, mixing, retry, and detector costs are paid. For
the stated safe stationary mass bound, the finite retry cap is logarithmic
in $1/\delta_{\rm prep}$, without a rare-event rejection cost.

The learned safe initial law has safe mass at least $e^{-2E_u}/2$ and
can therefore be sampled with finite capped rejection and prescribed
work. More generally a supplied safe-mass floor $s_*>0$ is possible,
but both physical preparation and safe-law sampling must pay
$O(s_*^{-1}\log(1/\delta_{\rm prep}))$ attempts (and the learned
density-ratio factor). The present constant-cost preparation claim uses
$s_*=1/2$ throughout. The density-ratio bound gives
\begin{equation}
\TV(\nu_u,\widehat\pi^u(\cdot\mid A))
\le \tanh(\beta_u/2)+\tanh(E_u).
\label{eq:prepTV}
\end{equation}
Add the numerical safe-sampling tolerance. Combining preparation,
Theorem~\ref{thm:stability} (or \eqref{eq:directStability}) and the finite
model sampler yields a completely specified error
\begin{equation}
\varepsilon_{{\rm law},u}=
\varepsilon_{{\rm model},u}+\varepsilon_{{\rm sim},u}
+\varepsilon_{{\rm init},u}
+\tanh(\beta_u/2)+\tanh(E_u).
\label{eq:totalAcquiredTV}
\end{equation}
This is a conditional-on-success mission law. For an unconditional
operation, output an explicit preparation-abort symbol and charge its
probability; assigning it loss at most one adds at most
$\delta_{\rm prep}$ to risk.

For every menu intervention, choose any bounded loss
$\ell_u(\mathcal R)\in[0,1]$. It may jointly depend on event time,
terminal damage and the retained precursor window. On the single
simultaneous acquisition/calibration confidence event, the difference
between physical and generated expected loss is at most
$\varepsilon_{{\rm law},u}$. With
\begin{equation}
n\ge\frac{\log(2K/\delta_{\rm dec})}{2\sigma^2}
\label{eq:decisionSamples}
\end{equation}
fresh generated records per intervention, all empirical risks have
additional error at most $\sigma$, with probability at least
$1-\delta_{\rm dec}$. Selecting their minimum gives
\begin{equation}
\operatorname{Risk}(\widehat u)
-\min_{u\in\mathcal U}\operatorname{Risk}(u)
\le2\left(\max_u\varepsilon_{{\rm law},u}+\sigma
+\delta_{\rm prep}\right).
\label{eq:regret}
\end{equation}
Any known intervention cost can be included after scaling total loss
into the declared bounded range. The event-law TV certificate covers
all bounded record losses, but the finite-sample simultaneous statement
above is for the declared $K$ evaluated losses; selecting arbitrary
losses after inspecting the same records needs its own complexity bound.

The decision is made after acquisition and before deploying the chosen
field and starting its preparation. Only the selected system is
physically prepared. This does not grant instantaneous feedback on an
unheld live network: observing a state, computing a choice and ramping
the apparatus take real time. A mission that must remain active during
acquisition or decision requires a separately priced holding, shutdown
or exposure model. The theorem does not supply such access.

Absolute-TV risk guarantees do not give relative accuracy for every
rare damage subtype, nor control unbounded lifetime costs. The acquired
relative whole-CDF theorem remains a separate operation and can be
used alongside the present joint-record guarantee.

